\chapter{Audio Features}

Chapter~26 turned continuous sound pressure into a stream of PCM samples, and Chapter~10
turned that stream into a spectrogram. Yet a raw waveform or spectrogram still contains
thousands of numbers per second. If you want a
machine to \emph{decide} something -- ``is this speech or silence?'', ``is this a
plucked string or a struck drum?'', ``did a note just start?'', ``which note is
it?'' -- you need to boil the rich, high-dimensional waveform down to a handful of
numbers that capture what matters and throw away what does not. Those numbers are
\textbf{features}.

A feature is a measurement. The waveform of one second of CD audio is
$44100$ numbers; that is far too much to feed a classifier directly, and almost
all of it is redundant or irrelevant for the question you actually care about. A
feature is a function that takes that wall of samples and returns something small
and meaningful -- the energy, the brightness, the pitch, the rate of change. Good
features are to a signal what a good summary is to a book: lossy on purpose,
keeping the gist.

This chapter builds the classic audio features from the ground up, geometrically.
We will see short-time energy and zero-crossing rate (the two cheapest features in
all of audio), the spectral shape descriptors (centroid, spread, rolloff, flux),
the mel scale and the MFCCs derived step by careful step, chroma for harmony,
pitch detection by autocorrelation and the YIN refinement, and onset detection for
rhythm. Throughout, keep one fact in mind: \textbf{every one of these features is
a function of either the waveform or its spectrum, and almost every one of them
applies, unchanged, to the motion signals from a wrist sensor.} Energy,
zero-crossing rate, spectral centroid, and flux are exactly the features a gesture
recognizer computes on accelerometer data. Learn them once here, on audio where
your ears can check the answer, and you get them for free on the IMU.

\chapterline

\section{What a feature is, geometrically}
\label{ch27:what-is-a-feature}

A frame of audio is a vector. Take $N$ consecutive samples
$\vec{x} = (x_0, x_1, \dots, x_{N-1})$ -- a typical frame is $20$--$40$\,ms, so
$N$ is a few hundred to a couple thousand samples. This frame lives in $\R^N$.
A feature is a map $\R^N \to \R$ (or to a small vector) that collapses those $N$
dimensions down to one or a few.

\begin{intuitionbox}
Think of the frame as a single point in a very high-dimensional space. A feature
is a \emph{measuring instrument} you point at that point. Energy measures how far
the point is from the origin. Zero-crossing rate measures how jagged the path
through the samples is. The spectral centroid measures, after rotating into the
Fourier basis, where the ``center of mass'' of the energy sits along the frequency
axis. Each feature is a different ruler, and the art is choosing rulers that line
up with the distinction you care about.
\end{intuitionbox}

Because audio changes over time, we almost never compute a feature on the whole
signal. We compute it \emph{per frame}, sliding a window along the signal exactly
as we did for the STFT in Chapter~10. The result is a \textbf{feature contour}: a
new, much shorter, much slower signal -- one number (or a few) per frame, perhaps
every $10$\,ms. We have traded $44100$ samples per second for, say, $100$ feature
vectors per second. This dimensional reduction is our goal: turn a dense audio stream into a
low-rate stream of descriptors that a downstream stage -- a threshold, a state
machine, a classifier -- can act on.

\begin{focusbox}
The pipeline for the entire chapter is always the same three steps:
\begin{enumerate}
  \item \textbf{Frame.} Slice the signal into overlapping windows (e.g. $25$\,ms
        frames, hopping $10$\,ms).
  \item \textbf{Window.} Multiply each frame by a taper (Hann, Hamming) to avoid
        edge artifacts -- see Chapter~10.
  \item \textbf{Measure.} Apply a feature function to each frame, producing one
        number (or vector) per frame.
\end{enumerate}
Different features differ only in step 3. Get the framing right once and you can
swap features freely.
\end{focusbox}

\begin{connectionbox}
A feature map $\R^N \to \R$ is what linear algebra (Chapter~42)
calls a \emph{functional}. The linear ones (like ``the DC value'' or any single
DFT bin) are inner products with a fixed vector; the nonlinear ones (energy is a
quadratic form, ZCR involves sign changes) are not, but the geometric picture --
projecting a high-dimensional point down to something measurable -- is the same.
And when a feature averages a quantity across a frame, it is a sample mean estimating an
\emph{expectation} (Chapter~20), so standard tools for means and variances apply
directly to feature contours.
\end{connectionbox}

\section{Short-time energy: how loud, right now}
\label{ch27:energy}

The simplest question you can ask of a frame is: how much is going on? The answer
is its \textbf{energy}, the squared length of the frame vector:
\begin{equation}
\label{ch27:energy-defn}
E[m] = \sum_{n=0}^{N-1} \bigl(x[mH + n]\, w[n]\bigr)^2 = \norm{\vec{x}_m}^2,
\end{equation}
where $m$ indexes the frame, $H$ is the hop size, and $w[n]$ is the window. This
is nothing but $\norm{\cdot}^2$ from Chapter~3, applied per frame. Silence sits
near the origin ($E \approx 0$); a loud passage sits far out.

Two cosmetic variants are common. The \textbf{root-mean-square (RMS)} amplitude
divides by $N$ and takes a square root, $\text{RMS}[m] = \sqrt{E[m]/N}$, which
puts it back in the same units as the samples (it is the standard deviation of the
frame when the mean is zero). And because loudness is perceived roughly
logarithmically, we usually report energy in \textbf{decibels},
$10\log_{10}(E[m] + \eps)$, with a tiny $\eps$ to avoid $\log 0$.

\begin{whybox}
Why square, rather than sum the absolute values? Squaring makes energy a smooth,
differentiable quadratic form $\vec{x}^\top\vec{x}$, it matches the physical
notion of energy (power is amplitude squared), and connects directly to
the spectrum through Parseval's theorem: the energy you compute in time
\eqref{ch27:energy-defn} equals (up to a constant) the sum of squared magnitudes
of the DFT. The same number, measured in two bases. That is why energy is the
bridge feature between the time-domain features of this section and the
spectral features of the next.
\end{whybox}

\begin{figure}[H]
\centering
\begin{tikzpicture}
\begin{axis}[
  width=12cm, height=3.6cm,
  xlabel={time (s)}, ylabel={$x(t)$},
  xmin=0, xmax=3, ymin=-1.2, ymax=1.2,
  axis lines=box, samples=200,
  title={A clip: quiet, then loud, then quiet},
  title style={font=\small},
]
\addplot[sigblue, thick, domain=0:1, samples=120] {0.15*sin(deg(2*pi*6*x))};
\addplot[sigblue, thick, domain=1:2, samples=160] {sin(deg(2*pi*6*x))};
\addplot[sigblue, thick, domain=2:3, samples=120] {0.15*sin(deg(2*pi*6*x))};
\end{axis}
\end{tikzpicture}

\vspace{0.3em}

\begin{tikzpicture}
\begin{axis}[
  width=12cm, height=3.2cm,
  xlabel={frame time (s)}, ylabel={RMS},
  xmin=0, xmax=3, ymin=0, ymax=0.85,
  axis lines=box,
  title={Short-time RMS contour tracks the loudness envelope},
  title style={font=\small},
]
\addplot[specred, thick] coordinates {
  (0,0.106)(0.5,0.106)(0.95,0.106)(1.05,0.707)(1.5,0.707)(1.95,0.707)
  (2.05,0.106)(2.5,0.106)(3,0.106)};
\end{axis}
\end{tikzpicture}
\caption{Top: a waveform that is quiet, then loud, then quiet. Bottom: the
short-time RMS contour, one value per frame, recovers the loudness envelope -- a
slow, low-rate signal you can threshold to find ``where the sound is.''}
\label{ch27:fig-energy}
\end{figure}

\begin{example}[Energy of a pure tone]
\label{ch27:ex-tone-energy}
Take a frame of a sinusoid $x[n] = A\sin(2\pi f n / f_s)$ with no window. Over an
integer number of cycles, $\sum \sin^2 = N/2$, so $E = A^2 N/2$ and
$\text{RMS} = A/\sqrt{2}$. That $\sqrt{2}$ is exactly the peak-to-RMS ratio of a
sine you may know from AC circuits: a sine of amplitude $1$ has RMS $0.707$. The
RMS contour of Figure~\ref{ch27:fig-energy} jumps from $0.15/\sqrt2 \approx 0.106$
to $1/\sqrt2 \approx 0.707$, which is precisely what the plot shows.
\end{example}

\begin{trapbox}
Energy is not loudness. Human loudness perception is frequency-dependent (a
$60$\,Hz tone and a $3$\,kHz tone of equal energy do \emph{not} sound equally
loud) and compressive. Raw RMS is a fine \emph{relative} measure for finding where
the action is, but do not call it loudness without an equal-loudness weighting
(A-weighting) and a logarithmic scale. The same caution applies to motion: the
acceleration ``energy'' of a flick and a slow wave can be similar while they feel
completely different.
\end{trapbox}

\section{Zero-crossing rate: how often it flips sign}
\label{ch27:zcr}

The second-cheapest feature counts how often the waveform crosses zero within a
frame:
\begin{equation}
\label{ch27:zcr-defn}
\text{ZCR}[m] = \frac{1}{2N}\sum_{n=1}^{N-1}
  \bigl|\,\operatorname{sgn}(x[mH+n]) - \operatorname{sgn}(x[mH+n-1])\,\bigr|.
\end{equation}
Each term is $1$ when consecutive samples have opposite signs and $0$ otherwise;
dividing by $2N$ gives a rate per sample (multiply by $f_s$ for crossings per
second). It costs only comparisons and additions -- no multiplies, no FFT -- which
is why early speech chips and the cheapest microcontrollers leaned on it heavily.

\begin{intuitionbox}
Zero-crossing rate is a crude pitch-and-noisiness meter rolled into one. A pure
low tone crosses zero twice per cycle, so its ZCR is small and proportional to
frequency. White noise wiggles constantly, so its ZCR is high. A periodic,
low-pitched sound (a vowel, a bowed cello) has low ZCR; a noisy, high-frequency
sound (the hiss of /s/, a cymbal, a snare) has high ZCR. So ZCR is the classic
\textbf{voiced/unvoiced} discriminator in speech, and more generally a
tonal-versus-noisy detector.
\end{intuitionbox}

\begin{figure}[H]
\centering
\begin{tikzpicture}
\begin{axis}[
  width=6cm, height=3.4cm,
  xlabel={time}, ylabel={},
  xlabel style={font=\footnotesize},
  xmin=0, xmax=2, ymin=-1.3, ymax=1.3,
  axis lines=left, xtick=\empty, ytick={-1,0,1},
  samples=120, domain=0:2,
  title={Low tone: low ZCR},
  title style={font=\small},
]
\draw[gray!50] (axis cs:0,0) -- (axis cs:2,0);
\addplot[sigblue, thick] {sin(deg(2*pi*1.5*x))};
\end{axis}
\end{tikzpicture}
\hspace{0.4cm}
\begin{tikzpicture}
\begin{axis}[
  width=6cm, height=3.4cm,
  xlabel={time}, ylabel={},
  xlabel style={font=\footnotesize},
  xmin=0, xmax=2, ymin=-1.3, ymax=1.3,
  axis lines=left, xtick=\empty, ytick={-1,0,1},
  samples=160, domain=0:2,
  title={Hissy/noisy: high ZCR},
  title style={font=\small},
]
\draw[gray!50] (axis cs:0,0) -- (axis cs:2,0);
\addplot[specred, thick] {0.7*sin(deg(2*pi*9*x))+0.3*sin(deg(2*pi*23*x))};
\end{axis}
\end{tikzpicture}
\caption{Zero-crossing rate distinguishes tonal from noisy. The slow tone (left)
flips sign rarely; the high-frequency, hissy signal (right) flips constantly.}
\label{ch27:fig-zcr}
\end{figure}

\begin{trapbox}
ZCR is fooled by DC offset and by noise near zero. A small constant bias can stop
a quiet signal from ever crossing zero (artificially low ZCR), while broadband
noise sitting on top of a clean tone adds spurious crossings each time the sum
dithers across zero (artificially high ZCR). Always remove the mean
(high-pass / DC-block) before computing ZCR, and consider a small hysteresis
band around zero so that tiny wiggles do not register as crossings.
\end{trapbox}

\begin{goalbox}
Short-time energy and ZCR are the first two features your wrist controller will
ever compute, and it will compute them on accelerometer magnitude, not audio.
Energy of the acceleration signal tells you \emph{how vigorous} a motion is --
the difference between resting and gesturing -- and is the natural trigger to wake
the recognizer from a low-power idle. ZCR of a gyroscope axis tells you how
\emph{oscillatory} the motion is -- a tremor or a fast shake has high ZCR; a
smooth sweep has low ZCR. Both cost a few additions per sample, which is exactly
what a battery-powered microcontroller on your forearm can afford to run
continuously. The caveat worth stating: these coarse features detect \emph{that} a
gesture happened and roughly how energetic it was; they do not, on their own,
recover which fingers moved. That finer question waits for Chapter~37.
\end{goalbox}

\section{Spectral shape: centroid, spread, rolloff}
\label{ch27:spectral-shape}

Now cross into the frequency domain. Take the magnitude spectrum of a frame,
$|X[k]|$ for bins $k = 0, \dots, N/2$ at frequencies $f_k = k f_s / N$. Treat the
\emph{energy per bin}, $|X[k]|^2$ (or sometimes $|X[k]|$), as a distribution of
``mass'' along the frequency axis. Once you see the spectrum as a mass
distribution, the language of statistics describes its shape exactly.

\subsection*{Spectral centroid: the center of mass (brightness)}

The \textbf{spectral centroid} is the energy-weighted mean frequency:
\begin{equation}
\label{ch27:centroid-defn}
\text{centroid}[m] = \frac{\sum_k f_k\, |X_m[k]|}{\sum_k |X_m[k]|}.
\end{equation}
This is precisely the formula for the center of mass of point masses $|X[k]|$ at
positions $f_k$ -- or, in probability language, the \emph{expected frequency}
$\E[f]$ under the distribution $p_k = |X[k]| / \sum_j |X[j]|$. A spectrum with most
of its energy down low (a bass note, a vowel) has a low centroid; one with energy
up high (a cymbal, the hiss of /s/) has a high centroid. The centroid correlates
remarkably well with the perceptual notion of \textbf{brightness}.

\begin{figure}[H]
\centering
\begin{tikzpicture}
\begin{axis}[
  width=12cm, height=4.2cm,
  xlabel={frequency (kHz)}, ylabel={$|X[k]|$},
  xmin=0, xmax=8, ymin=0, ymax=1.55,
  axis lines=box, ytick=\empty,
  title={Two spectra and their centroids},
  title style={font=\small},
]
\addplot[sigblue, thick, smooth] coordinates {
  (0,0.1)(0.5,0.6)(1,1.0)(1.5,0.8)(2,0.45)(3,0.2)(4,0.08)(6,0.03)(8,0.01)};
\addplot[specred, thick, smooth] coordinates {
  (0,0.02)(1,0.05)(2,0.15)(3,0.4)(4,0.8)(5,1.0)(6,0.7)(7,0.4)(8,0.2)};
\draw[sigblue, dashed, thick] (axis cs:1.25,0) -- (axis cs:1.25,1.1);
\draw[specred, dashed, thick] (axis cs:4.7,0) -- (axis cs:4.7,1.1);
\node[sigblue] at (axis cs:1.5,1.28) {\small low centroid (dark)};
\node[specred] at (axis cs:5.6,1.28) {\small high centroid (bright)};
\end{axis}
\end{tikzpicture}
\caption{The spectral centroid is the balance point of the magnitude spectrum.
The blue spectrum (energy low) balances around $1.25$\,kHz; the red spectrum
(energy high) balances around $4.7$\,kHz and sounds brighter.}
\label{ch27:fig-centroid}
\end{figure}

\subsection*{Spectral spread, skewness, flatness}

If the centroid is the mean, the \textbf{spectral spread} (or bandwidth) is the
standard deviation -- how wide the energy is smeared around the centroid:
\begin{equation}
\label{ch27:spread-defn}
\text{spread}[m] = \sqrt{\frac{\sum_k (f_k - \text{centroid}[m])^2\,|X_m[k]|}
                              {\sum_k |X_m[k]|}}.
\end{equation}
A pure tone has near-zero spread; broadband noise has huge spread. Carrying the
moment-by-moment analogy further gives spectral \emph{skewness} (lopsidedness) and
\emph{kurtosis} (peakiness), all borrowed verbatim from the moments of a
probability distribution.

A close cousin is \textbf{spectral flatness}, the ratio of the geometric mean to
the arithmetic mean of the spectrum:
\begin{equation}
\label{ch27:flatness-defn}
\text{flatness}[m] =
  \frac{\bigl(\prod_k |X_m[k]|^2\bigr)^{1/K}}
       {\frac{1}{K}\sum_k |X_m[k]|^2}.
\end{equation}
Flatness is near $1$ for white noise (a flat spectrum: geometric mean equals
arithmetic mean) and near $0$ for a pure tone (energy in one bin makes the
geometric mean collapse). It is the standard \textbf{tonal-versus-noisy} measure,
the spectral-domain partner of ZCR, and it is what audio codecs use to decide how
aggressively they can quantize a band.

\subsection*{Spectral rolloff: the frequency below which most energy lives}

The \textbf{spectral rolloff} $R[m]$ is the frequency below which a fixed fraction
$\kappa$ (commonly $0.85$ or $0.95$) of the total spectral energy is contained:
\begin{equation}
\label{ch27:rolloff-defn}
R[m] = \min\Bigl\{ f_K : \sum_{k \le K} |X_m[k]|^2
                          \ge \kappa \sum_{k} |X_m[k]|^2 \Bigr\}.
\end{equation}
In probability terms this is just the $\kappa$-\emph{quantile} of the spectral
energy distribution. Rolloff, like centroid, tracks brightness and the
voiced/unvoiced distinction, but it is more robust to a few stray
high-frequency bins because it asks about cumulative energy rather than a mean.

\begin{example}[Centroid of a two-tone frame]
\label{ch27:ex-centroid}
Suppose a frame's spectrum has just two significant lines: amplitude $3$ at
$200$\,Hz and amplitude $1$ at $1000$\,Hz. The centroid is
\[
\frac{3\cdot 200 + 1\cdot 1000}{3 + 1} = \frac{1600}{4} = 400\text{ Hz}.
\]
The strong low line pulls the balance point down toward it -- exactly as a heavier
mass pulls a center of gravity. If you added high-frequency noise the centroid
would climb, which is the basis for using it as a brightness/noise detector.
\end{example}

\section{Spectral flux: how fast the spectrum is changing}
\label{ch27:flux}

Everything so far described a \emph{single} frame. \textbf{Spectral flux} measures
the \emph{change} between consecutive frames -- how different this frame's spectrum
is from the last one:
\begin{equation}
\label{ch27:flux-defn}
\text{flux}[m] = \sum_k \bigl(\, |X_m[k]| - |X_{m-1}[k]| \,\bigr)_{+}^2,
\end{equation}
where $(\cdot)_+ = \max(\cdot, 0)$ keeps only the bins that \emph{grew}
(half-wave rectification). Geometrically, flux is (a rectified version of) the
squared Euclidean distance between two successive spectrum vectors: it is large
when the sound's frequency content suddenly shifts and near zero when the sound is
steady.

\begin{whybox}
Why keep only the increases? Because we usually use flux to find \emph{onsets} --
moments when new energy appears (a note begins, a consonant bursts, a drum is
struck). When a note \emph{ends}, energy in its bins \emph{decreases}; that is a
release, not an onset, and we do not want it to trigger. Half-wave rectification
keeps only rising energy, so flux peaks at the births of sounds, not their deaths.
This makes spectral flux the core tool of onset detection in
Section~\ref{ch27:onset}.
\end{whybox}

\begin{figure}[H]
\centering
\begin{tikzpicture}
\begin{axis}[
  width=12cm, height=3.4cm,
  xlabel={frame time (s)}, ylabel={flux},
  xmin=0, xmax=4, ymin=0, ymax=1.1,
  axis lines=box, ytick=\empty,
  title={Spectral flux spikes when new notes begin},
  title style={font=\small},
]
\addplot[structurepurple, thick] coordinates {
  (0,0.02)(0.48,0.03)(0.5,0.95)(0.55,0.2)(0.6,0.04)(0.9,0.03)
  (1.0,0.85)(1.05,0.18)(1.1,0.03)(1.7,0.02)(1.8,0.9)(1.85,0.22)(1.9,0.03)
  (2.5,0.02)(2.6,0.8)(2.65,0.15)(2.7,0.03)(3.3,0.02)(3.4,0.92)(3.45,0.2)(3.5,0.03)(4,0.02)};
\end{axis}
\end{tikzpicture}
\caption{Spectral flux over a clip of five separated notes. Each note onset
produces a sharp spike; the steady-state and silence between notes produce almost
no flux. Picking the peaks of this curve detects onsets.}
\label{ch27:fig-flux}
\end{figure}

\begin{connectionbox}
Flux is a discrete time-derivative of the spectrum, a finite difference quotient (Chapter~42)
applied here to a vector-valued
signal (the whole spectrum at each frame). Many features in this chapter are
``derivative-like'': flux is the derivative of the spectrum, onset strength is the
derivative of energy, and the famous \emph{delta} and \emph{delta-delta} MFCCs are
literally first and second derivatives of the cepstral features over time. A
recurring theme: \emph{what} something is comes from one frame; \emph{events}
come from how it changes between frames.
\end{connectionbox}

\section{The mel scale: hearing as the ear does}
\label{ch27:mel}

So far our spectral features used a \emph{linear} frequency axis: bin $k$ sits at
$f_k = k f_s/N$, equally spaced. But human hearing is not linear. We are
exquisitely sensitive to changes among low frequencies and progressively coarser
up high. The musical interval from $100$ to $200$\,Hz (one octave) sounds like the
same ``distance'' as $1000$ to $2000$\,Hz, even though the second spans ten times
as many hertz. Pitch perception is roughly \emph{logarithmic}.

The \textbf{mel scale} is an empirical warping of the frequency axis built to make
equal perceptual distances correspond to equal numeric distances. A common formula
is
\begin{equation}
\label{ch27:mel-formula}
m(f) = 2595 \, \log_{10}\!\Bigl(1 + \frac{f}{700}\Bigr),
\end{equation}
which is nearly linear below about $1$\,kHz and logarithmic above it -- matching
the ear's behavior. ``Mel'' comes from ``melody''; the scale was anchored so that
$1000$\,mel equals a $1000$\,Hz tone at a reference loudness.

\begin{figure}[H]
\centering
\begin{tikzpicture}
\begin{axis}[
  width=10cm, height=5.2cm,
  xlabel={frequency $f$ (Hz)}, ylabel={mel $m(f)$},
  xmin=0, xmax=8000, ymin=0, ymax=3000,
  axis lines=box, samples=120, domain=0:8000,
  title={The mel scale warps frequency: fine resolution low, coarse high},
  title style={font=\small},
]
\addplot[sigblue, thick] {2595*log10(1+x/700)};
\addplot[gray, dashed] {x*2521/8000};
\end{axis}
\end{tikzpicture}
\caption{The mel mapping (blue) versus a straight line (gray). Equal steps on the
vertical mel axis correspond to ever-wider hertz steps as frequency rises, which
is why a mel-spaced filterbank places many narrow filters down low and a few wide
ones up high.}
\label{ch27:fig-mel}
\end{figure}

\subsection*{The mel filterbank}

We do not resample the spectrum onto the mel axis directly; instead we pool the
linear spectrum into a small number of overlapping \textbf{triangular filters}
spaced evenly on the mel axis. Each filter is a triangle in frequency: it rises
from zero at the previous center, peaks at its own center, and falls to zero at
the next center. Because the centers are equally spaced \emph{in mel}, the
triangles are narrow at low frequency and wide at high frequency. We compute, for
each of (typically) $26$--$40$ filters, the weighted sum of spectral energy
falling under its triangle. The output is the \textbf{mel spectrum}: a short
vector that summarizes the spectral envelope the way the ear's frequency analysis
roughly does.

\begin{figure}[H]
\centering
\begin{tikzpicture}
\begin{axis}[
  width=12.5cm, height=4.6cm,
  title={Mel filterbank: overlapping triangular filters with psychoacoustic spacing},
  title style={font=\small\bfseries},
  xlabel={frequency $f$ (Hz)}, ylabel={filter weight $H_m(f)$},
  xmin=0, xmax=8000, ymin=0, ymax=1.28,
  axis lines=box, ytick={0,0.5,1.0},
  xtick={0,1000,2000,3000,4000,5000,6000,7000,8000},
  xticklabels={0,1k,2k,3k,4k,5k,6k,7k,8k},
]
% Regime labels
\draw[gray!60, dashed] (axis cs:1000, 0) -- (axis cs:1000, 1.25);
\node[sourceblue, font=\scriptsize\bfseries] at (axis cs:500, 1.15) {Linear regime ($f < 1$\,kHz)};
\node[specred!80!black, font=\scriptsize\bfseries] at (axis cs:4500, 1.15) {Logarithmic regime ($f > 1$\,kHz)};

% Narrow low-frequency triangles
\addplot[sigblue, thick, fill=fillsig, fill opacity=0.35] coordinates {(0,0)(150,1)(330,0)} \closedcycle;
\addplot[sigblue, thick, fill=fillsig, fill opacity=0.35] coordinates {(150,0)(330,1)(560,0)} \closedcycle;
\addplot[sigblue, thick, fill=fillsig, fill opacity=0.35] coordinates {(330,0)(560,1)(850,0)} \closedcycle;
\addplot[sigblue, thick, fill=fillsig, fill opacity=0.35] coordinates {(560,0)(850,1)(1230,0)} \closedcycle;

% Mid-frequency triangles
\addplot[specred, thick, fill=specred, fill opacity=0.25] coordinates {(850,0)(1230,1)(1750,0)} \closedcycle;
\addplot[specred, thick, fill=specred, fill opacity=0.25] coordinates {(1230,0)(1750,1)(2450,0)} \closedcycle;
\addplot[specred, thick, fill=specred, fill opacity=0.25] coordinates {(1750,0)(2450,1)(3400,0)} \closedcycle;

% High-frequency wide triangles
\addplot[structurepurple, thick, fill=structurepurple, fill opacity=0.25] coordinates {(2450,0)(3400,1)(4700,0)} \closedcycle;
\addplot[structurepurple, thick, fill=structurepurple, fill opacity=0.25] coordinates {(3400,0)(4700,1)(6400,0)} \closedcycle;
\addplot[structurepurple, thick, fill=structurepurple, fill opacity=0.25] coordinates {(4700,0)(6400,1)(8000,0)} \closedcycle;
\end{axis}
\end{tikzpicture}
\caption{A mel filterbank with overlapping triangular weighting windows. Filter peak frequencies are equidistant on the perceptual mel scale. When mapped back to hertz, filters are narrow and densely packed at low frequencies (where human pitch discrimination is finest) and grow progressively wider at higher frequencies (mimicking cochlear critical bands).}
\label{ch27:fig-melbank}
\end{figure}

\begin{historybox}
The foundations of modern audio features united psychophysics, geophysics, and speech
engineering. In 1937, Harvard psychologists Stanley Smith Stevens, John Volkmann, and
Edwin Newman asked human listeners to adjust audio tones until one pitch sounded ``twice
as high'' as another, establishing the empirical mel scale. In 1963, Bruce Bogert,
M.~Healy, and John Tukey coined the term \emph{cepstrum} (an anagram of spectrum) while
analyzing seismic echoes from underground nuclear tests, realizing that the logarithm of
a spectrum turned multiplicative echo convolutions into separable additions. In 1980,
Steven Davis and Paul Mermelstein united these ideas into Mel-Frequency Cepstral
Coefficients (MFCCs). The innovation sparked debate: traditional phoneticians
argued that speech recognition required explicit physical models of vocal-tract formants
and poles, rather than ``black-box'' empirical cepstral projections, while hearing
scientists debated whether the mel scale or physiological Equivalent Rectangular
Bandwidth (ERB) filters better matched the cochlea. Davis and Mermelstein's approach
won on empirical performance, becoming the dominant acoustic front end for four decades
across speech recognition, speaker verification, and music classification. An active open
challenge is designing learned, end-to-end neural audio front ends that match the compact
interpretability of MFCCs while remaining robust to severe reverberation and acoustic domain shift.
\end{historybox}

\section{MFCCs, derived step by step}
\label{ch27:mfcc}

\textbf{Mel-frequency cepstral coefficients} (MFCCs) are the single most famous
audio feature -- the backbone of speech recognition for decades and still a strong
baseline for any audio classification task. They look intimidating as a formula
but each step has a clear geometric purpose. Let us build them from the mel
spectrum we just computed.

\begin{focusbox}
The MFCC recipe, in seven steps:
\begin{enumerate}
  \item \textbf{Frame and window} the signal (Hann/Hamming, $\sim$25\,ms).
  \item \textbf{FFT}; take the magnitude (or power) spectrum $|X[k]|^2$.
  \item \textbf{Mel filterbank}: pool into $\sim$26--40 mel-band energies.
  \item \textbf{Logarithm}: take $\log$ of each band energy.
  \item \textbf{DCT}: apply a discrete cosine transform across the bands.
  \item \textbf{Keep the low coefficients} (typically the first 12--13).
  \item \textbf{(Optional) deltas}: append first/second time-derivatives.
\end{enumerate}
\end{focusbox}

Steps 1--3 we have already met. The new leverage comes in steps 4 and 5; here is why each
one is there.

\subsection*{Step 4: the logarithm separates source from filter}

Speech (and most pitched instruments) is well modeled as a \emph{source} -- the
buzzing vocal folds or vibrating string, rich in harmonics -- passed through a
\emph{filter} -- the resonances of the vocal tract or instrument body that shape
which harmonics are loud. In the frequency domain these \emph{multiply}: the
observed spectrum is (source spectrum) $\times$ (filter response). Multiplication
is hard to separate; but $\log(ab) = \log a + \log b$ turns the product into a
\emph{sum}. After the log, the slowly varying filter envelope and the rapidly
varying harmonic source sit side by side as additive components, ready to be
pried apart.

\begin{happeningbox}
This is the central trick of the \emph{cepstrum} (an anagram of ``spectrum'').
Think of the log-spectrum as itself a kind of signal, but one whose ``time'' axis
is frequency. The vocal-tract filter is a \emph{slowly} varying shape across
frequency -- a few broad bumps (the formants). The harmonic source is a
\emph{fast} ripple across frequency -- a comb of evenly spaced harmonics. ``Slow''
and ``fast'' variation across the frequency axis means \emph{low} and \emph{high}
``frequency'' of this new signal. So a Fourier-like transform of the log-spectrum
sorts filter (slow, low) from source (fast, high). Keeping the low coefficients
keeps the filter shape -- the timbre -- and discards the pitch. That is exactly
what you want for recognizing \emph{what} was said regardless of \emph{what pitch}
it was said at.
\end{happeningbox}

\begin{recoverybox}
\textbf{Pause and think}: If the vocal tract filter multiplies the vocal fold excitation spectrum, why can we not separate them by division? In the linear frequency domain, $X(f) = E(f) H(f)$, so deconvolution by direct spectral division ($X(f)/E(f)$) blows up whenever $E(f)$ is near zero between harmonics. Taking the logarithm converts multiplication into addition, $\log |X| = \log |E| + \log |H|$, allowing smooth low-quefrency liftering to isolate the vocal tract envelope $\log |H|$ without dividing by zero.
\end{recoverybox}

\subsection*{Step 5: the DCT decorrelates and compresses}

To Fourier-analyze the log-mel spectrum we use the \textbf{discrete cosine
transform} (DCT) rather than the full DFT. The DCT expresses the $M$-vector of
log-mel energies as a sum of slow-to-fast cosine shapes:
\begin{equation}
\label{ch27:dct}
c[i] = \sum_{j=0}^{M-1} \log\!\bigl(S[j]\bigr)\,
       \cos\!\Bigl[\frac{\pi i}{M}\Bigl(j + \tfrac12\Bigr)\Bigr],
       \qquad i = 0, 1, \dots, M-1,
\end{equation}
where $S[j]$ is the $j$-th mel band energy. The coefficients $c[i]$ are the
MFCCs. We keep the low-order ones ($i = 1, \dots, 12$ or so), which describe the
broad, smooth shape of the spectral envelope, and drop the high-order ones, which
capture the fine harmonic ripple (pitch) we want to ignore.

\begin{whybox}
Why the DCT and not the DFT? Three reasons, all geometric. First, the log-mel
spectrum is real, and the DCT is the natural real-input cousin of the DFT (it
assumes an even extension, so it needs only cosines -- no imaginary part, no phase
to worry about). Second, the DCT is famously good at \emph{energy compaction}: for
the smooth, correlated log-mel envelopes that speech produces, it packs almost all
the information into the first few coefficients, so $12$ numbers suffice. Third --
and this is the deep reason -- for signals with the kind of correlation structure
speech has, the DCT basis is very nearly the Karhunen--Lo\`eve (eigenvector)
basis, the optimal one that \emph{decorrelates} the components. Decorrelated
features are exactly what simple classifiers (which assume diagonal covariance)
want. The DCT is doing a free, fixed, fast approximation to PCA.
\end{whybox}

\subsection*{What the coefficients mean}

$c[0]$ is the average of the log-mel energies -- essentially overall log-energy
(loudness), and often dropped or replaced by the frame energy. $c[1]$ measures the
overall \emph{tilt} of the spectrum (energy concentrated low versus high -- a
brightness/balance axis, close cousin of the centroid). $c[2]$ measures whether
energy is concentrated in the middle versus the edges, and so on: each higher
coefficient corresponds to a faster cosine ripple across the mel bands, describing
ever-finer detail of the spectral envelope shape.

\begin{figure}[H]
\centering
\begin{tikzpicture}[font=\small, >=Stealth]
\tikzset{
  stg/.style={draw=blue!70!black, fill=sourceblue!10, rounded corners=2pt, minimum height=10mm, minimum width=15mm, align=center, thick},
  mel/.style={draw=targetgreen!70!black, fill=targetgreen!10, rounded corners=2pt, minimum height=10mm, minimum width=16mm, align=center, thick},
  cep/.style={draw=quotientorange!80!black, fill=quotientorange!12, rounded corners=2pt, minimum height=10mm, minimum width=16mm, align=center, thick}
}

\node[stg] (a) at (0,0) {Frame \&\\Window};
\node[stg] (b) at (2.4,0) {FFT\\$|X[k]|^2$};
\node[mel] (c) at (4.9,0) {Mel Filter\\Bank};
\node[mel] (d) at (7.1,0) {Logarithm\\$\log(E_m)$};
\node[cep] (e) at (9.3,0) {DCT\\Type-II};
\node[cep] (f) at (11.7,0) {Truncate\\$c_1 \dots c_{12}$};

\draw[->, thick] (a) -- (b);
\draw[->, thick] (b) -- (c);
\draw[->, thick] (c) -- (d);
\draw[->, thick] (d) -- (e);
\draw[->, thick] (e) -- (f);

% Subtle functional bracket labels below
\draw[gray!50, thick] (0,-0.8) -- (2.4,-0.8) node[midway, below, font=\tiny\scshape, text=gray!80] {Time $\to$ Freq};
\draw[gray!50, thick] (4.9,-0.8) -- (7.1,-0.8) node[midway, below, font=\tiny\scshape, text=gray!80] {Auditory Scale};
\draw[gray!50, thick] (9.3,-0.8) -- (11.7,-0.8) node[midway, below, font=\tiny\scshape, text=gray!80] {Decorrelation};
\end{tikzpicture}
\caption{The complete MFCC extraction pipeline. Framing and FFT compute the raw power spectrum; the mel filterbank pools energy into critical bands; the logarithm converts multiplicative vocal-tract filtering into additive components; and the DCT decorrelates the bands, compacting the spectral envelope into 12--13 cepstral coefficients.}
\label{ch27:fig-mfcc-pipeline}
\end{figure}

\begin{example}[Why the log matters: a worked sketch]
\label{ch27:ex-log}
Imagine two recordings of the same vowel /a/, one sung at $110$\,Hz and one at
$220$\,Hz. The vocal-tract filter (the formant shape that makes it an /a/) is
identical; only the harmonic comb differs. Before the log, doubling the source
amplitude or changing its harmonic spacing scrambles the spectrum
multiplicatively. After $\log$, the formant envelope is the same additive
``baseline'' in both, and the harmonic combs differ only in their fast ripple --
which the DCT shoves into high coefficients you discard. The surviving low MFCCs
are nearly identical for both pitches: the feature has become
\emph{pitch-invariant}, recognizing the /a/-ness. That invariance is the entire
reason MFCCs dominated speech recognition.
\end{example}

\begin{trapbox}
MFCCs deliberately throw away pitch. That is a feature for speech recognition and
a \emph{bug} for tasks where pitch is the point -- melody transcription, key
detection, singer identification. For those, use chroma (next section) or a pitch
tracker. Choosing MFCCs for a melody task is one of the most common beginner
mistakes in music information retrieval.
\end{trapbox}

\subsection*{Deltas: motion in the cepstrum}

A single MFCC vector describes one $25$\,ms snapshot, but speech and gesture are
about \emph{trajectories}. The \textbf{delta} coefficients approximate the time
derivative of each MFCC,
\begin{equation}
\label{ch27:delta}
\Delta c[i,m] \approx \frac{\sum_{\theta=1}^{\Theta}\theta\,(c[i,m+\theta]-c[i,m-\theta])}
                            {2\sum_{\theta=1}^{\Theta}\theta^2},
\end{equation}
a regression slope over a few neighboring frames, and the \textbf{delta-delta}
coefficients do it again for the second derivative. Appending deltas typically
triples the feature dimension ($13 \to 39$) and captures \emph{how the spectrum is
moving} -- which, as we keep seeing, is where the events live.

\section{Chroma: folding the spectrum onto the twelve pitch classes}
\label{ch27:chroma}

For music we often want the opposite of MFCCs: we want to keep harmony and pitch
class while ignoring timbre and octave. The \textbf{chroma} (or pitch-class
profile) feature does this by folding the spectrum onto the twelve notes of the
chromatic scale -- C, C\#, D, \dots, B -- summing the energy of all octaves of
each pitch class into a single bin.

The construction uses the same logarithmic ear-fact as the mel scale, but applied
musically: an octave is a doubling of frequency, and Western music splits each
octave into twelve equal ratios of $2^{1/12}$ (the equal-tempered semitone). To
build chroma, map each spectral bin to its nearest musical pitch, then to that
pitch's class modulo $12$ (collapsing octaves), and accumulate energy into a
$12$-dimensional vector:
\begin{equation}
\label{ch27:chroma-defn}
\text{chroma}[p] = \sum_{k\,:\;\text{class}(f_k)=p} |X[k]|^2,
\qquad p = 0,\dots,11.
\end{equation}

\begin{figure}[H]
\centering
\begin{tikzpicture}
\begin{axis}[
  width=12cm, height=4cm,
  xlabel={pitch class}, ylabel={energy},
  xmin=-0.6, xmax=11.6, ymin=0, ymax=1.1,
  axis lines=box, ytick=\empty,
  xtick={0,1,2,3,4,5,6,7,8,9,10,11},
  xticklabels={C,C\#,D,D\#,E,F,F\#,G,G\#,A,A\#,B},
  title={Chroma vector of a C-major chord (C--E--G)},
  title style={font=\small},
]
\addplot[ybar, bar width=10pt, fill=targetgreen, draw=targetgreen!60!black] coordinates {
  (0,1.0)(1,0.08)(2,0.1)(3,0.06)(4,0.9)(5,0.07)(6,0.05)(7,0.85)
  (8,0.06)(9,0.09)(10,0.05)(11,0.07)};
\end{axis}
\end{tikzpicture}
\caption{Chroma of a C-major chord: energy piles up on C, E, and G regardless of
which octaves the notes are played in or what instrument plays them. The small
nonzero values on other classes are leakage from harmonics and windowing.}
\label{ch27:fig-chroma}
\end{figure}

\begin{intuitionbox}
Chroma is the musical inverse of MFCCs. MFCCs keep the spectral \emph{envelope}
(timbre) and discard pitch; chroma keeps the \emph{pitch class} and discards both
timbre and octave. A C-major chord on a piano and on a guitar have very different
MFCCs (different timbres) but nearly identical chroma (same notes). That octave-
and timbre-invariance is exactly what you want for chord recognition, key
estimation, and matching cover versions of a song.
\end{intuitionbox}

\begin{applicationbox}
Chroma underlies practical music systems: chord-recognition apps that print
guitar tabs from audio, key-detection in DJ software that matches songs in
compatible keys, and audio-to-score alignment that follows a live performer
through a known piece. Comparing the sequence of chroma vectors between two
recordings (via dynamic time warping) is also how ``cover song'' and
audio-fingerprint-adjacent systems decide that two performances are the same
underlying song.
\end{applicationbox}

\section{Pitch detection: finding the fundamental period}
\label{ch27:pitch}

Pitch is not the same as a single feature value -- it is an estimate of the
\textbf{fundamental frequency} $f_0$, the rate at which a (quasi-)periodic signal
repeats. A vowel, a bowed string, a sung note: all repeat a basic waveform shape
many times per frame. Pitch detection asks: what is the period of that repetition?

\begin{trapbox}
``Just take the FFT and find the biggest peak.'' This fails constantly. For many
instruments and voices the fundamental is \emph{weaker} than its harmonics, or
even physically absent (the ``missing fundamental'' -- a telephone removes
everything below $300$\,Hz, yet you still hear a $150$\,Hz male voice at its true
pitch because your brain infers $f_0$ from the harmonic \emph{spacing}). The
strongest spectral peak is often the second or third harmonic. Robust pitch
detection looks at \emph{periodicity}, not at peak height.
\end{trapbox}

\subsection*{Autocorrelation: how well a signal matches a shifted copy of itself}

The cleanest way to find a period is to slide the signal against a copy of itself
and see at which shift it lines up best. That is the \textbf{autocorrelation}:
\begin{equation}
\label{ch27:autocorr}
r[\tau] = \sum_{n} x[n]\,x[n+\tau].
\end{equation}
At lag $\tau = 0$ the signal perfectly overlaps itself, giving the maximum (that
is just the energy). As $\tau$ grows, $r[\tau]$ dips and then, when $\tau$ equals
one period $T_0$, the signal lines up with itself again and $r[\tau]$ peaks once
more. The first strong peak after lag zero sits at $\tau = T_0$, and
$f_0 = f_s / T_0$.

\begin{whybox}
Autocorrelation is an inner product (Chapter~3) between the signal and a
time-shifted version of itself: $r[\tau] = \inner{x[n]}{x[n+\tau]}$. An inner
product measures alignment -- how much two vectors point the same way. When the
shift equals the period, the shifted copy is nearly identical to the original, so
their alignment (inner product) is large. Pitch detection by autocorrelation is
therefore literally ``at what shift does the signal most resemble itself?'' -- a
purely geometric question.
\end{whybox}

\begin{figure}[H]
\centering
\begin{tikzpicture}
\begin{axis}[
  name=ax1,
  width=12cm, height=3.2cm,
  xlabel={$n$}, ylabel={$x[n]$},
  xmin=0, xmax=4, ymin=-1.6, ymax=1.6,
  axis lines=left, samples=160, domain=0:4,
  xtick={0,1,2,3,4}, ytick={-1,0,1},
  title={A quasi-periodic frame},
  title style={font=\small},
]
\draw[gray!40] (axis cs:0,0) -- (axis cs:4,0);
\addplot[sigblue, thick] {sin(deg(2*pi*x)) + 0.4*sin(deg(2*pi*2*x)) + 0.2*sin(deg(2*pi*3*x))};
\end{axis}
\begin{axis}[
  at={(ax1.south west)}, anchor=north west, yshift=-1.9cm, clip=false,
  width=12cm, height=3.2cm,
  xlabel={lag $\tau$}, ylabel={$r[\tau]$},
  xmin=0, xmax=4, ymin=-1.1, ymax=1.2,
  axis lines=left, samples=160, domain=0:4,
  xtick={0,1,2,3,4}, ytick={-1,0,1},
  title={Its autocorrelation: peak at lag zero, next peak at the period $T_0$},
  title style={font=\small},
]
\draw[gray!40] (axis cs:0,0) -- (axis cs:4,0);
\addplot[specred, thick] {(cos(deg(2*pi*x)) + 0.16*cos(deg(2*pi*2*x)) + 0.04*cos(deg(2*pi*3*x)))/1.2};
\draw[specred!80, dashed, thick] (axis cs:1,-1.1) -- (axis cs:1,1.0);
\node[anchor=south west, font=\small, specred] at (axis cs:1.05,0.85) {$\tau = T_0$};
\end{axis}
\end{tikzpicture}
\caption{Top: a periodic frame. Bottom: its autocorrelation peaks at lag $0$ and
again at the period $T_0$. Locating that first non-trivial peak gives the period,
hence $f_0 = f_s/T_0$.}
\label{ch27:fig-autocorr}
\end{figure}

\begin{connectionbox}
By the convolution theorem (Chapter~11), autocorrelation can be computed
\emph{fast} as $r = \text{IFFT}\bigl(|X[k]|^2\bigr)$: the inverse transform of the
power spectrum. So pitch detection in time (autocorrelation) and pitch analysis in
frequency (the power spectrum) are two views of the same object -- the
Wiener--Khinchin theorem made practical. This also reveals why ``peak in spectrum''
and ``peak in autocorrelation'' are related but not identical: the autocorrelation
folds all the harmonics back onto the single period, which is exactly why it sees
the true $f_0$ even when the fundamental bin is weak.
\end{connectionbox}

\subsection*{YIN: fixing autocorrelation's mistakes}

Plain autocorrelation has two well-known failure modes: it can pick lag zero or a
too-short lag (\emph{octave-too-high} errors, since a half-period also correlates
somewhat), and amplitude changes within the frame bias the peaks. The
\textbf{YIN} algorithm (de Cheveign\'e and Kawahara, 2002) is a celebrated set of
refinements that makes period detection far more reliable. Its key ideas:

\begin{enumerate}
  \item \textbf{Difference function instead of correlation.} Use the squared
        difference $d[\tau] = \sum_n (x[n] - x[n+\tau])^2$, which is \emph{zero}
        when the shifted copy matches (a valley to find) rather than a peak. This
        is more stable than correlation under amplitude variation.
  \item \textbf{Cumulative mean normalization.} Divide $d[\tau]$ by its running
        average up to $\tau$. This deliberately makes lag zero \emph{not} a winner
        and suppresses the spurious sub-period dips that cause octave errors.
  \item \textbf{Absolute threshold.} Take the first valley that drops below a fixed
        threshold (e.g. $0.1$), rather than the global minimum -- this prefers the
        true (longest, lowest-pitch) period over its harmonics.
  \item \textbf{Parabolic interpolation.} Fit a parabola around the chosen valley
        to estimate the period to sub-sample precision, beating the resolution of
        the integer lags.
\end{enumerate}

YIN and its descendants (pYIN, CREPE in the neural era) are what real pitch
trackers use. The lesson is general: the naive idea (autocorrelation) is correct
in spirit, and the engineering is in defeating its predictable biases -- octave
errors, amplitude sensitivity, and discrete resolution.

\begin{example}[From lag to frequency]
\label{ch27:ex-lag}
At $f_s = 16{,}000$\,Hz, suppose the difference function's first deep valley sits
at lag $\tau = 145$ samples, and parabolic interpolation refines it to
$\tau = 145.3$. Then
\[
f_0 = \frac{f_s}{\tau} = \frac{16000}{145.3} \approx 110.1\text{ Hz},
\]
an A2 -- the open A string of a bass guitar. Note how the interpolation matters:
the integer lag $145$ would give $110.34$\,Hz and lag $146$ would give
$109.6$\,Hz, a difference you can hear. Sub-sample interpolation is not a luxury
for pitch.
\end{example}

\section{Onset detection: when did something happen}
\label{ch27:onset}

Rhythm is about \emph{events in time}: the instants when notes begin, drums hit,
syllables start. An \textbf{onset} is the start of a new sonic event. Detecting
onsets turns continuous audio into a sequence of time stamps -- the raw material
for beat tracking, tempo estimation, automatic transcription, and rhythm games.

The standard recipe builds an \textbf{onset strength signal} (also called a
detection function), one value per frame that is large at onsets and small
otherwise, then picks its peaks. The building blocks are features we already have:

\begin{itemize}
  \item \textbf{Energy difference.} The simplest detector: the positive
        first difference of the (log) energy contour. Works for percussive,
        sharply-attacked sounds (drums, plucks) where energy jumps suddenly.
  \item \textbf{Spectral flux.} Section~\ref{ch27:flux}, the half-wave-rectified
        change in the magnitude spectrum. Catches onsets that change \emph{timbre}
        or pitch without a big energy jump (a legato note change, a bowed
        attack) -- the most widely used general-purpose onset feature.
  \item \textbf{Phase deviation.} For soft, blurry onsets, the \emph{phase} of
        each bin advances predictably during steady tones; an onset disrupts that
        predictability. Phase-based and complex-domain detectors combine magnitude
        flux with phase irregularity for the hardest cases.
\end{itemize}

\begin{figure}[H]
\centering
\begin{tikzpicture}
\begin{axis}[
  width=12cm, height=3.6cm,
  xlabel={time (s)}, ylabel={onset strength},
  xmin=0, xmax=4, ymin=0, ymax=1.2,
  axis lines=box, ytick=\empty,
  title={Onset strength with peak-picking threshold},
  title style={font=\small},
]
\addplot[structurepurple, thick] coordinates {
  (0,0.05)(0.3,0.07)(0.4,0.95)(0.5,0.3)(0.6,0.08)(1.0,0.06)(1.1,0.8)(1.2,0.25)
  (1.3,0.07)(1.55,0.5)(1.6,0.9)(1.7,0.3)(1.8,0.08)(2.4,0.05)(2.5,0.85)(2.6,0.28)
  (2.7,0.07)(3.2,0.06)(3.3,0.92)(3.4,0.3)(3.5,0.08)(4,0.05)};
\addplot[trapred, dashed, thick] coordinates {(0,0.45)(4,0.45)};
\node[trapred] at (axis cs:3.4,0.58) {\small threshold};
\end{axis}
\end{tikzpicture}
\caption{An onset-strength contour. Peaks above an adaptive threshold (red) are
declared onsets. The threshold usually tracks a local moving average so that quiet
and loud passages are handled fairly.}
\label{ch27:fig-onset}
\end{figure}

\subsection*{Peak picking}

Turning the strength curve into discrete onsets needs care, because the curve is
noisy and onsets come in clusters. Practical peak picking: (1) smooth the strength
signal lightly; (2) require a peak to exceed a local moving-average threshold by a
margin, so loud and quiet sections are treated fairly; (3) enforce a minimum
spacing between onsets (a refractory period of a few tens of milliseconds) so a
single attack is not counted several times. These three steps -- smooth,
adaptive-threshold, refractory -- recur in every event detector, audio or not.

\begin{applicationbox}
In the 2000s, systematic benchmark evaluations across shared musical corpora (such as
the MIREX campaigns) compared spectral flux, complex-domain, and phase-based detectors.
The broad finding -- that spectral flux provides an exceptionally solid, simple
baseline for percussive attacks, while soft onsets in expressive bowed strings or vocals
benefit from tracking phase regularity -- established the standard event-detection
pipeline used in digital audio workstations and realtime gesture controllers alike.
\end{applicationbox}

\section{Putting it together: a feature pipeline on a clip}
\label{ch27:pipeline}

The features in this chapter are rarely used alone; a real system computes several
per frame and stacks them into a feature vector, then feeds the stream to a
classifier or a state machine. Here is a compact, correct sketch of the front end
in Python/NumPy.

\begin{lstlisting}[language=Python, caption={Framing plus a handful of features per frame.}]
import numpy as np

def frames(x, N=1024, H=512):
    # generator of windowed frames
    w = np.hanning(N)
    for m in range(0, len(x) - N, H):
        yield x[m:m+N] * w

def features(frame, fs):
    # time-domain
    energy = np.sum(frame**2)
    rms    = np.sqrt(energy / len(frame))
    sgn    = np.sign(frame)
    zcr    = np.mean(np.abs(np.diff(sgn))) / 2.0
    # spectral
    X  = np.abs(np.fft.rfft(frame))
    fk = np.fft.rfftfreq(len(frame), 1.0 / fs)
    Xs = X.sum() + 1e-12
    centroid = np.sum(fk * X) / Xs
    spread   = np.sqrt(np.sum((fk - centroid)**2 * X) / Xs)
    cumE     = np.cumsum(X**2)
    rolloff  = fk[np.searchsorted(cumE, 0.85 * cumE[-1])]
    return dict(rms=rms, zcr=zcr, centroid=centroid,
                spread=spread, rolloff=rolloff)

def spectral_flux(prev_mag, mag):
    d = mag - prev_mag
    return np.sum(np.maximum(d, 0.0)**2)   # half-wave rectified
\end{lstlisting}

\begin{walkbox}
Run this front end on a clip and you get five contours (RMS, ZCR, centroid,
spread, rolloff) plus a flux contour, each at the frame rate. Look at them
together: during a vowel, RMS is high, ZCR is low, centroid is low; during an /s/,
RMS is moderate, ZCR is high, centroid is high. During silence everything drops.
You have, in a few dozen lines, built a representation in which speech-versus-
silence and voiced-versus-unvoiced are nearly linearly separable -- the kind of
input a tiny classifier can handle. Ask yourself: which single feature best
separates your particular clip's classes? That is the feature-selection question
every applied project must answer.
\end{walkbox}

\begin{buildbox}
On a microcontroller you will not run an FFT every frame for free, so budget
deliberately. The cheap tier -- RMS energy and ZCR -- needs no FFT and runs
comfortably at full sample rate on an Arduino-class chip; use it as an
always-on gate. The spectral tier -- centroid, rolloff, flux -- needs one
real FFT per frame; a $256$- or $512$-point fixed-point FFT (e.g. CMSIS-DSP's
\texttt{arm\_rfft}) on a Cortex-M4 at $80$--$120$\,MHz fits easily in a $10$\,ms
hop. The expensive tier -- a full mel filterbank plus DCT for MFCCs -- is feasible
on an M4/M7 but is where you start watching the clock and the battery. A common
embedded pattern: gate on energy, and only run the FFT/MFCC stage when the gate is
open. This same tiered budgeting is exactly how the wrist controller will manage
its own feature extraction.
\end{buildbox}

\section{The same features on motion signals}
\label{ch27:motion}

Here is the payoff promised at the start. A tri-axial accelerometer streams three
signals $a_x[n]$, $a_y[n]$, $a_z[n]$ at, say, $100$--$200$\,Hz; a gyroscope adds three
more. These are signals just like audio, only slower and lower-dimensional, and
the features of this chapter transfer with almost no change:

\begin{itemize}
  \item \textbf{Energy / RMS} of the acceleration magnitude
        $\sqrt{a_x^2+a_y^2+a_z^2}$ measures motion intensity -- resting versus
        gesturing -- and is the standard activity gate.
  \item \textbf{Zero-crossing rate} of a band-passed axis measures how oscillatory
        a motion is, distinguishing a shake or tremor from a smooth sweep.
  \item \textbf{Spectral centroid and rolloff} of a short motion FFT capture
        whether the motion energy is slow (a deliberate wave) or fast (a flick or
        vibration).
  \item \textbf{Spectral flux} marks the \emph{onset} of a gesture, the motion
        analogue of a note attack -- the trigger that says ``a gesture is
        starting, begin classifying.''
  \item \textbf{MFCC-style features} (a small mel-like filterbank plus DCT on the
        motion spectrum) are used in human-activity-recognition research and work
        surprisingly well for distinguishing walking, typing, and gesturing.
\end{itemize}

\begin{figure}[H]
\centering
\begin{tikzpicture}
\begin{axis}[
  width=12cm, height=3.4cm,
  xlabel={time (s)}, ylabel={accel mag (g)},
  xmin=0, xmax=4, ymin=0, ymax=2.2,
  axis lines=box, samples=200, domain=0:4,
  title={Accelerometer magnitude: rest, gesture, rest},
  title style={font=\small},
]
\addplot[sigblue, thick, domain=0:1.2, samples=60] {1 + 0.03*sin(deg(2*pi*3*x))};
\addplot[specred, thick, domain=1.2:2.6, samples=120]
  {1 + 0.8*sin(deg(2*pi*4*x))*exp(-1.2*(x-1.9)^2)};
\addplot[sigblue, thick, domain=2.6:4, samples=60] {1 + 0.03*sin(deg(2*pi*3*x))};
\end{axis}
\end{tikzpicture}
\caption{A motion trace with a gesture in the middle. Short-time energy rises in
the gesture region (intensity gate); spectral flux spikes at its start (onset);
centroid reports how fast the motion is. The very same code from
Section~\ref{ch27:pipeline} produces these contours.}
\label{ch27:fig-motion}
\end{figure}

\begin{goalbox}
For the wrist controller, the plan is now concrete. An always-on energy gate on
accelerometer magnitude wakes the system; a spectral-flux onset detector marks the
gesture's start; and a short feature vector (energy, ZCR, centroid, a few
mel-DCT coefficients, plus their deltas) computed over the gesture window feeds the
classifier that decides \emph{which} gesture it was, which is then mapped to a MIDI
message. Every block in that sentence is a feature from this chapter.

\textbf{The limit, stated plainly:} this pipeline recognizes coarse,
discrete wrist and forearm gestures -- a flick, a twist, a tap, a hold. A
wrist-worn IMU \emph{cannot} recover continuous individual finger joint angles;
the wrist simply does not move enough, distinctly enough, when only a finger
moves. Detecting which finger pressed, or tracking continuous finger pose, needs
sensors that see the muscles or the fingers themselves -- surface EMG or force
myography on the forearm for discrete finger events, or instrumented rings for
continuous pose. We confront those sensors and their own features in
Chapter~37. Audio-style features get you robust \emph{gesture} control; they do
not get you a glove.
\end{goalbox}

\section{Summary}
\label{ch27:summary}

A feature collapses a high-dimensional frame to a few meaningful numbers, computed
frame by frame to produce a low-rate contour. The cheapest features are
time-domain: \emph{energy/RMS} (squared length, the loudness envelope) and
\emph{zero-crossing rate} (sign flips, a tonal-versus-noisy and voiced-versus-
unvoiced meter). Treating the magnitude spectrum as a mass distribution yields the
shape features: \emph{centroid} (mean frequency, brightness), \emph{spread}
(standard deviation), \emph{rolloff} (a quantile), and \emph{flatness}
(tonal-versus-noisy). \emph{Spectral flux} measures change between frames and
drives onset detection.

The \emph{mel scale} warps frequency to match the ear's logarithmic pitch
sensitivity; pooling the spectrum through a mel filterbank, taking the log
(to separate source from filter), and applying a DCT (to decorrelate and compress)
yields the \emph{MFCCs}, the timbre-describing, pitch-invariant feature that
dominated speech recognition. Their opposite, \emph{chroma}, folds the spectrum
onto twelve pitch classes to keep harmony while discarding timbre and octave.
\emph{Pitch} is found by periodicity -- autocorrelation (an inner product of the
signal with a shifted copy of itself), refined by \emph{YIN} to beat octave errors
and reach sub-sample precision. \emph{Onsets} come from peak-picking a detection
function built from energy difference, spectral flux, or phase deviation.

Every one of these features is a function of the waveform or its spectrum, and
nearly all of them apply unchanged to motion signals, which is why this chapter is
the feature toolbox for the wrist controller -- with the standing caveat that IMU
features yield coarse gestures, not finger pose.

With a rich vocabulary for analyzing incoming sounds and motion signals, we turn
in the next chapter from analysis to creation: audio effects and synthesis. We will see
how delay lines, comb filters, wavetables, and physical models generate living sounds,
and how the features we built here can modulate sound parameters in real time.

\section{Exercises}
\label{ch27:exercises}

\begin{exercise}[RMS and dB]
\label{ch27:ex1}
A frame of $N=1024$ samples holds a sine of amplitude $0.5$ over an integer number
of cycles. Compute its energy, RMS, and energy in dB ($10\log_{10}E$). Then a
second frame holds the same sine plus white noise of RMS $0.05$; estimate the new
RMS (assume signal and noise are uncorrelated, so energies add) and the change in
dB. Is the dB change something a listener would notice?
\end{exercise}

\begin{exercise}[ZCR of a tone]
\label{ch27:ex2}
Show that a pure sinusoid of frequency $f$ sampled at $f_s$ has an expected
zero-crossing rate of $2f/f_s$ crossings per sample. Use this to ``read off'' the
frequency of a clean tone from its ZCR. Then explain, with a sketch, why this
estimate breaks down when a second, higher-frequency tone or noise is added.
\end{exercise}

\begin{exercise}[Centroid arithmetic]
\label{ch27:ex3}
A frame's spectrum has lines of magnitude $4$ at $250$\,Hz, $2$ at $500$\,Hz, and
$1$ at $2000$\,Hz. Compute the spectral centroid and the spectral spread. Now
double the $2000$\,Hz line to magnitude $2$ and recompute both. By how much, and in
which direction, did each move? Interpret in terms of brightness.
\end{exercise}

\begin{exercise}[Rolloff as a quantile]
\label{ch27:ex4}
Given the same three-line spectrum as Exercise~\ref{ch27:ex3} (use magnitudes
$4, 2, 1$ as energies for simplicity), compute the $0.85$ spectral rolloff
frequency. Explain why rolloff can be more robust than centroid when a single
loud high-frequency bin appears, and give an example spectrum where centroid moves
a lot but rolloff barely moves.
\end{exercise}

\begin{exercise}[Why the log in MFCCs]
\label{ch27:ex5}
Two frames have spectra $A[k]$ and $2A[k]$ (the second is the first scaled by a
constant gain, as if recorded louder). Show that after the $\log$ step their
mel-band vectors differ only by a constant offset, and therefore -- because the
DCT of a constant lands entirely in $c[0]$ -- their MFCCs $c[1], c[2], \dots$ are
\emph{identical}. What practical invariance does this give the feature?
\end{exercise}

\begin{exercise}[DCT energy compaction]
\label{ch27:ex6}
Take a length-8 vector that is a smooth ramp $v_j = j$ for $j=0,\dots,7$. Compute
its DCT (by hand or with a few lines of code) and report how much of the total
energy lands in the first two coefficients versus the rest. Repeat for a vector of
random values. Explain what this demonstrates about why MFCCs keep only the
low-order coefficients.
\end{exercise}

\begin{exercise}[Autocorrelation period]
\label{ch27:ex7}
A signal is sampled at $f_s = 8000$\,Hz. Its autocorrelation has its first strong
peak after lag zero at $\tau = 100$ samples. What is the estimated $f_0$? Now
suppose a spurious peak also appears at $\tau = 50$; which note is that, and why
might a naive ``first peak'' rule pick it? Which YIN step is designed to prevent
this error?
\end{exercise}

\begin{exercise}[The missing fundamental]
\label{ch27:ex8}
A periodic tone has $f_0 = 200$\,Hz but the recording high-passes away everything
below $400$\,Hz, leaving only the harmonics at $400, 600, 800, \dots$\,Hz. Explain
why an FFT-peak pitch detector would report the wrong answer, while an
autocorrelation detector still finds $f_0 = 200$\,Hz. What is the lag of the
relevant autocorrelation peak at $f_s = 16000$\,Hz?
\end{exercise}

\begin{exercise}[Onset detection function]
\label{ch27:ex9}
Explain why spectral flux uses half-wave rectification (keeping only increases)
rather than the full difference. Construct a short example -- a note that grows then
fades -- and sketch both the full-difference and the half-wave-rectified flux,
showing how the rectified version peaks at the attack only.
\end{exercise}

\begin{exercise}[Peak picking pitfalls]
\label{ch27:ex10}
You apply a fixed global threshold to an onset-strength curve and find it misses
soft onsets in quiet passages while double-counting loud drum hits. Propose two
fixes (one about the threshold, one about timing), and explain how each addresses a
specific failure. Relate your answer to the smooth / adaptive-threshold /
refractory recipe.
\end{exercise}

\begin{exercise}[Transfer to motion]
\label{ch27:ex11}
You record $100$\,Hz accelerometer magnitude during ten repetitions of a wrist
flick separated by rest. Describe which features from this chapter you would
compute and how you would combine an energy gate with a spectral-flux onset
detector to segment the ten flicks. Then state plainly what this pipeline
\emph{cannot} tell you about the hand, and which sensor you would add to recover
finger-level information.
\end{exercise}

\begin{exercise}[Feature separability]
\label{ch27:ex12}
For each pair below, name the single feature from this chapter that best separates
them and justify it in one sentence: (a) silence vs.\ speech; (b) a vowel vs.\ the
consonant /s/; (c) a flute note vs.\ a cymbal crash; (d) a C-major chord vs.\ an
A-minor chord; (e) a sustained note vs.\ the instant a new note begins.
\end{exercise}
